\subsection{The compact detector's finite ordinary normalizer comparison}

\begin{result}{A.4.22}{Proved}{Exact compact ordinary reconstruction and the final inner quotient}
Let $E=\mathcal P_c$ be the order-$32\cdot3^{14}$ target
of A.4.18.  Retain its characteristic subgroup
$A\cong C_3^{12}$ and quotient $B=G_2\times H$ from
(EX.1).  On the faithful quaternion multiplicity space
let $\mathscr N$ be the finite group of pairs
\[
 (\bar g,q),\qquad
 \bar g= g|_{M/\mathfrak m_H^2M},\quad
 g\in N_{\operatorname{SO}_{18}^-(3)}(H),\quad
 q\in\Gamma(2)\bmod12,
                                                        \tag{NC.1}
\]
with the two exact compatibilities
\[
 \nu_g=q\bmod3,
 \qquad \theta(g)=\eta(q),\qquad
 \eta\!\left(I+2\begin{pmatrix}\alpha&\beta\\
                         \gamma&\alpha\end{pmatrix}\right)
       =(-1)^{\alpha+\beta+\gamma}.                     \tag{NC.2}
\]
The first equality uses the original downstairs deck
labels $a^4,b^4$.  In the primitive upstairs handle
basis $(A=b^3,B=a^3)$ the period matrix is conjugated
by the retained swap $J$; therefore this equality
is $\nu_g=Jq_{\rm up}J\pmod3$ in that basis.
The character $\eta$ is unchanged by this swap.
The spinor in (NC.2) is well-defined on $\bar g$:
the kernel of reduction to $M/\mathfrak m_H^2M$
is a finite three-group, so its spinor norm is one.

There is an actual homomorphism
\[
 \mathscr N\longrightarrow\operatorname{Out}(E)
                                                        \tag{NC.3}
\]
obtained from (EX.2)--(EX.5), with exact kernel the
specified inner $Z(Q_8)\times H$ of order $18$.
Its image is precisely the image of the ordinary
$\Gamma(2)\leq\operatorname{SL}_2(\mathbb Z)$ on $E$.
The exact orders are
\[
 |\mathscr N|=384\cdot3^6,
 \qquad |\operatorname{im}(\Gamma(2)\to
                   \operatorname{Out}(E))|=192\cdot3^4.
                                                        \tag{NC.4}
\]
After restoring the six ordinary modulo-two classes,
the whole ordinary modular image on $E$ has order
$1152\cdot3^4$.  Its previously computed detected
quotient $J_c$ is therefore an isomorphic image;
its CM centralizer census in A.4.18 applies to the
actual ordinary outer surface image itself.
\end{result}

\begin{proof}
The independent hermitian normalizer calculation
in A.4.21 proves that
$\mathscr N$ has at most $384\cdot3^6$ elements.
We recall the exact count to specify its inputs.
The group algebra is $\mathbb F_3[H]$, and its
augmentation quotient by $\mathfrak m_H^2$ has basis
$1,X,Y$ with $X^2=XY=Y^2=0$ and involution
$\bar X=-X$, $\bar Y=-Y$.  The reduced hermitian
form on the two free module generators is
anisotropic.  The constant matrix of a special
orthogonal semilinear normalizer is in
$\operatorname{SO}_2^-(3)=C_4$.
For each fixed label map and constant matrix,
the two first-order coefficients satisfy two
self-adjoint $2\times2$ matrix conditions and
therefore have at most $3^{3+3}$ possibilities.
The spinor condition selects at most two of the
four constants.  There are $24$ possible label
maps $q\bmod3$ and eight possible $q\bmod4$.
This gives the displayed bound, with all its
period, label and spinor constraints retained.

\textbf{Actual source inclusion and the matching order.}
The program
\path{work/tower_code/quaternion_compact_normalizer_joint.py}
starts from the original four free-surface maps
$u,l,\operatorname{Ad}(a),\operatorname{Ad}(b)$.
Their quaternion-automorphism residues generate
all $24$ elements.  The complete $24\cdot4=96$
Schreier words of the kernel are retained, together
with their actual signed period matrices.
The program restricts them to the twelve-dimensional
image of $T$, and then to its faithful primitive
multiplicity factor of dimension six.  A specified
free basis over $\mathbb F_3[H]/\mathfrak m_H^2$
identifies the latter with two constant coordinates,
two $X$ coordinates and two $Y$ coordinates.
It checks, for every one of the $96$ words, the
literal conjugation identities
\[
 g h_jg^{-1}=h_1^{q_{1j}}h_2^{q_{2j}}\qquad(j=1,2).
                                                        \tag{NC.5}
\]
Thus the direction of the label/period map is
checked, not inferred from matching orders.

The separate bridge producer
\path{work/tower_code/quaternion_compact_normalizer_bridge.py}
rebuilds the same four ordinary surface actions on
all $73$ free-cover homology generators.  It gives
an explicit rank-six matrix $P:M_{18}\to M_6$ and
verifies all $96$ identities
\[
                         P g_{18}=g_6P.                \tag{NC.6}
\]
The original free words, the $18\times18$ matrices,
and $P$ are retained.  The independent audit checks
their special orthogonality and their spinor/period
identities.  These are ordinary geometric maps,
so they lie in the specified group (NC.1)--(NC.2).

In these coordinates the semisimple data are the
constant $2\times2$ matrix and the signed period
matrix.  The actual generated base has $384$
elements.  All $384\cdot96=36864$ Cayley-Schreier
rows are computed in the eight possible first-order
coordinates.  Their span has dimension six and is
invariant under all $96$ actual source generators.
Schreier generation proves that this is the entire
kernel.  The actual ordinary modular/deck image
therefore has order $384\cdot3^6$.
It is a subgroup of $\mathscr N$ of the upper-bound
order, so it equals $\mathscr N$.
The independent reproduction rebuilds the source
words and both parts of this enumeration, and
checks that all six kernel directions satisfy
the hermitian self-adjoint equations.

\textbf{The marked extension and its exact inner kernel.}
The module map represented by $\bar g$ acts on
$A=U_2\otimes M_6$ as $I\otimes\bar g$.
It centralizes the specified quaternion action
and has the $H$-label action in (NC.5).
Together with the literal base map (EX.5), this
satisfies (EX.2), so it induces an automorphism
of $E$ by (EX.3).  Changing the chosen splitting
changes it only by its specified inner-$A$
conjugation.  Hence the outer map (NC.3) is
well-defined and is a homomorphism.

An inner automorphism of $E$ has identity period
matrix.  In the quaternion-centralizing sector,
its base element must lie over $Z(Q_8)\times H$:
in $G_2$, the preimage of $Z(Q_8)$ is exactly
$Z(G_2)$, as follows by taking commutators in the
literal fibre product (EX.1).
The additional central kernel of $G_2\to Q_8$
acts trivially on $A$.  The remaining actual
inner action on $A$ is therefore precisely
$Z(Q_8)\times H$, of order $18$; it is faithful
since $-1$ acts as $-I$ and the nine translation
matrices are distinct.  Conversely every one
of these $18$ actions is induced by an actual
inner element of the original $\Sigma$ and
has identity period.  These observations prove
both inclusions for the kernel of (NC.3).

The quaternion-automorphism kernel of the four
ordinary generators maps onto $\Gamma(2)$.
Indeed an ordinary matrix is identity modulo two
exactly when its induced quaternion automorphism
is inner; the actual conjugations by $a,b$
realize all quaternion inner automorphisms and
remove that residue without changing the
ordinary outer matrix.  Therefore the image of
(NC.3), computed from the actual source words,
is exactly the ordinary $\Gamma(2)$ image.
Dividing the proved source order by $18$ gives
(NC.4).

Finally the signed period projection of the full
ordinary modular image retains its full
$\operatorname{SL}_2(\mathbb F_2)$ quotient of
order six.  The image of its kernel is the
$\Gamma(2)$ image just computed.  The full
ordinary image consequently has order
$6\cdot192\cdot3^4=1152\cdot3^4$.
The surjective detected-image map to $J_c$ has
the same finite order by A.4.18, so it is an
isomorphism.  This proves the final assertion
without assuming a hidden restriction kernel
vanishes.
\end{proof}

\paragraph{Scope of the remaining finite fusion.}
The proposition proves the exact ordinary comparison
for the finite group defined by (NC.1)--(NC.2).
In the nine-point application the actual joint
peripheral quotient must still provide the displayed
label/period equality, and the actual fusion must
preserve that lower quotient while putting the
faithful translation matrices into this specified
normalizer.  The earlier class-three peripheral
comparison and the special-orthogonal correction
are the available source maps for that step;
neither requirement is discarded in (NC.1).
